% =====================================================================
%  比較：標準の手法との比較
%  文の元：paper/sections/6-comparison.tex．表は論文の表 3（\NUM{tab-*}，\NUM{cmp-*}）
%  図の値：compare_k0.csv，compare_k5.csv，kcurves.csv，entmin.csv・entmin_mean.csv，finetune_k0.csv・finetune_k5.csv
% =====================================================================
\setband{比較　標準の手法と比べる}

\begin{frame}{9 つの標準の手法を，同じ生徒・同じ出発点・同じ選び方で比べる}
  \rl{比較の相手}
  \vspace{0.1em}
  \centering\tiny
  \renewcommand{\arraystretch}{1.15}
  \begin{tabularx}{\textwidth}{@{}p{0.8cm} p{4.4cm} X p{4.0cm}@{}}
    \toprule
    略称 & 手法 & すること & 渡す先生 \\
    \midrule
    ST & 自己訓練（Lee 2013） & 生徒自身の自信のある予測をラベルにする & -- \\
    HPL & ハードな擬似ラベル & 先生の予測のうちしきい値を超えたものを硬いラベルにする & 融合の出発点 \\
    EntM & エントロピー最小化（Grandvalet \& Bengio 2004） & ラベル無しの発話で予測のエントロピーを下げる & -- \\
    KD & 温度つきの蒸留（Hinton ら 2015） & 温度で軟らかくした先生の確率で学ぶ & 融合の出発点，または単独の出発点の積 \\
    GD & 一般化蒸留（Lopez-Paz ら 2016） & 特権情報の先生から蒸留する & $K$ 本で学び直した融合の生徒 \\
    CoT & 共訓練（Blum \& Mitchell 1998） & モダリティごとの分類器が互いにラベルを付ける & -- \\
    CoR & 共正則化（Sindhwani ら 2005） & モダリティごとの予測が一致するように罰する & -- \\
    PA & 事前確率の調整（Saerens ら 2002） & EM で感情の事前確率を推定し，出発点の出力を直す & -- \\
    RB & 出力の釣り合い直し & 出発点の出力を感情の間で釣り合わせ直す & -- \\
    \bottomrule
  \end{tabularx}\par
  \vspace{0.3em}
  {\raggedright\tiny KD の先生は，単独の生徒には融合の出発点，融合の生徒には単独の出発点の積．融合の生徒では GD と HPL は定義できず，CoT は 2 つの変種で入る\par}
  \vspace{0.2em}
  \takeaway{拡張版で使う「生徒の積の先生」はどの比較手法にも渡さない．差には先生の選び方も含まれる}
\end{frame}

\begin{frame}{比べ方：再学習する手法はすべて同じ縛りつきの生徒で，設定は仮のターゲットで選ぶ}
  \rl{比べ方}
  \vspace{0.3em}
  \begin{itemize}\footnotesize\setlength{\itemsep}{7pt}
    \item 分類器を再学習する手法は，すべて同じ縛りつきの線形の生徒を，同じ 5 つの値を選んだ出発点から，同じ $K$ 本と同じプールで学ぶ．事前確率の調整と釣り合い直しは出発点の出力を直す
    \item \textbf{調整版}：しきい値・温度・ラベル無しの損失の比重を，私たちの手法と同じ仮のターゲットで $K$ ごとに選ぶ．\textbf{固定版}：元の論文の値を使う
    \item 基本版は第 II 部の基本の手法，拡張版はその候補を足した手法
    \item 評価する生徒：AVT の世界の音声・顔・言葉・融合の生徒（ターゲット \NUM{cmp-targets-avt-n} つ），AV・VT・AT の融合の生徒（\NUM{cmp-targets-av-n}・\NUM{cmp-targets-vt-n}・\NUM{cmp-targets-at-n} つ），全部のエンコーダを 2 つ目のもの（HuBERT，CLIP，DeBERTa）に替えた AVT の 4 つの生徒．計 11 の生徒
    \item ターゲットごとに反復 \NUM{cmp-reps-n} 回の平均．仮のターゲットは反復 \NUM{cmp-pseudo-reps-n} 回
  \end{itemize}
\end{frame}

\begin{frame}{拡張版は $K=0$・5 のすべての生徒で，最も強い比較手法を上回る（論文の表 3）}
  \rl{結果：比較の表}
  \vspace{-0.5em}
  \centering\scriptsize
  \setlength{\tabcolsep}{4.2pt}
  \renewcommand{\arraystretch}{0.98}
  \begin{tabular}{@{}llrrlrrrrrr@{}}
    \toprule
     & & \multicolumn{2}{c}{出発点（UAR）} & 最も強い & \multicolumn{2}{c}{\textcolor{cBest}{最も強い − 出発点}} & \multicolumn{2}{c}{\textcolor{cOurs}{基本版 − 出発点}} & \multicolumn{2}{c}{\textcolor{cExt}{拡張版 − 出発点}} \\
    \cmidrule(lr){3-4}\cmidrule(lr){6-7}\cmidrule(lr){8-9}\cmidrule(l){10-11}
    世界 & 生徒 & $K{=}0$ & $K{=}5$ & （$K{=}0$/5） & $K{=}0$ & $K{=}5$ & $K{=}0$ & $K{=}5$ & $K{=}0$ & $K{=}5$ \\
    \midrule
    AVT & A & \NUM{tab-avt-a-k0-e39} & \NUM{tab-avt-a-k5-e39} & GD/GD & \NUM{cmp-best-vs-e39-avt-a-k0} & \NUM{cmp-best-vs-e39-avt-a-k5} & \NUM{cmp-ours-vs-e39-avt-a-k0} & \NUM{cmp-ours-vs-e39-avt-a-k5} & \NUM{cmp-oursx-vs-e39-avt-a-k0} & \NUM{cmp-oursx-vs-e39-avt-a-k5} \\
     & V & \NUM{tab-avt-v-k0-e39} & \NUM{tab-avt-v-k5-e39} & RB/なし & \NUM{cmp-best-vs-e39-avt-v-k0} & \NUM{cmp-best-vs-e39-avt-v-k5} & \NUM{cmp-ours-vs-e39-avt-v-k0} & \NUM{cmp-ours-vs-e39-avt-v-k5} & \NUM{cmp-oursx-vs-e39-avt-v-k0} & \NUM{cmp-oursx-vs-e39-avt-v-k5} \\
     & T & \NUM{tab-avt-t-k0-e39} & \NUM{tab-avt-t-k5-e39} & HPL/GD & \NUM{cmp-best-vs-e39-avt-t-k0} & \NUM{cmp-best-vs-e39-avt-t-k5} & \NUM{cmp-ours-vs-e39-avt-t-k0} & \NUM{cmp-ours-vs-e39-avt-t-k5} & \NUM{cmp-oursx-vs-e39-avt-t-k0} & \NUM{cmp-oursx-vs-e39-avt-t-k5} \\
     & 融合 & \NUM{tab-avt-f-k0-e39} & \NUM{tab-avt-f-k5-e39} & RB/ST & \NUM{cmp-best-vs-e39-avt-f-k0} & \NUM{cmp-best-vs-e39-avt-f-k5} & \NUM{cmp-ours-vs-e39-avt-f-k0} & \NUM{cmp-ours-vs-e39-avt-f-k5} & \NUM{cmp-oursx-vs-e39-avt-f-k0} & \NUM{cmp-oursx-vs-e39-avt-f-k5} \\
    AV & 融合 & \NUM{tab-av-f-k0-e39} & \NUM{tab-av-f-k5-e39} & RB/ST & \NUM{cmp-best-vs-e39-av-f-k0} & \NUM{cmp-best-vs-e39-av-f-k5} & \NUM{cmp-ours-vs-e39-av-f-k0} & \NUM{cmp-ours-vs-e39-av-f-k5} & \NUM{cmp-oursx-vs-e39-av-f-k0} & \NUM{cmp-oursx-vs-e39-av-f-k5} \\
    VT & 融合 & \NUM{tab-vt-f-k0-e39} & \NUM{tab-vt-f-k5-e39} & RB/ST & \NUM{cmp-best-vs-e39-vt-f-k0} & \NUM{cmp-best-vs-e39-vt-f-k5} & \NUM{cmp-ours-vs-e39-vt-f-k0} & \NUM{cmp-ours-vs-e39-vt-f-k5} & \NUM{cmp-oursx-vs-e39-vt-f-k0} & \NUM{cmp-oursx-vs-e39-vt-f-k5} \\
    AT & 融合 & \NUM{tab-at-f-k0-e39} & \NUM{tab-at-f-k5-e39} & RB/CoT & \NUM{cmp-best-vs-e39-at-f-k0} & \NUM{cmp-best-vs-e39-at-f-k5} & \NUM{cmp-ours-vs-e39-at-f-k0} & \NUM{cmp-ours-vs-e39-at-f-k5} & \NUM{cmp-oursx-vs-e39-at-f-k0} & \NUM{cmp-oursx-vs-e39-at-f-k5} \\
    \midrule
    AVT & A & \NUM{tab-w2-a-k0-e39} & \NUM{tab-w2-a-k5-e39} & CoT/GD & \NUM{cmp-best-vs-e39-w2-a-k0} & \NUM{cmp-best-vs-e39-w2-a-k5} & \NUM{cmp-ours-vs-e39-w2-a-k0} & \NUM{cmp-ours-vs-e39-w2-a-k5} & \NUM{cmp-oursx-vs-e39-w2-a-k0} & \NUM{cmp-oursx-vs-e39-w2-a-k5} \\
    （2 つ目の & V & \NUM{tab-w2-v-k0-e39} & \NUM{tab-w2-v-k5-e39} & CoT/KD & \NUM{cmp-best-vs-e39-w2-v-k0} & \NUM{cmp-best-vs-e39-w2-v-k5} & \NUM{cmp-ours-vs-e39-w2-v-k0} & \NUM{cmp-ours-vs-e39-w2-v-k5} & \NUM{cmp-oursx-vs-e39-w2-v-k0} & \NUM{cmp-oursx-vs-e39-w2-v-k5} \\
    \quad エンコーダ） & T & \NUM{tab-w2-t-k0-e39} & \NUM{tab-w2-t-k5-e39} & RB/GD & \NUM{cmp-best-vs-e39-w2-t-k0} & \NUM{cmp-best-vs-e39-w2-t-k5} & \NUM{cmp-ours-vs-e39-w2-t-k0} & \NUM{cmp-ours-vs-e39-w2-t-k5} & \NUM{cmp-oursx-vs-e39-w2-t-k0} & \NUM{cmp-oursx-vs-e39-w2-t-k5} \\
     & 融合 & \NUM{tab-w2-f-k0-e39} & \NUM{tab-w2-f-k5-e39} & RB/CoT & \NUM{cmp-best-vs-e39-w2-f-k0} & \NUM{cmp-best-vs-e39-w2-f-k5} & \NUM{cmp-ours-vs-e39-w2-f-k0} & \NUM{cmp-ours-vs-e39-w2-f-k5} & \NUM{cmp-oursx-vs-e39-w2-f-k0} & \NUM{cmp-oursx-vs-e39-w2-f-k5} \\
    \bottomrule
  \end{tabular}\par
  \vspace{0.2em}
  {\raggedright\tiny 差は UAR のポイントで，ターゲットと反復 10 回の平均．「なし」は，調整版の比較手法がどれもラベル無しの発話を使わない設定を選んだこと．略称は前のスライド\par}
  \vspace{0.15em}
  \takeaway{拡張版は全部で \NUM{cmp-oursx-vs-best-range-lo}〜\NUM{cmp-oursx-vs-best-range-hi} \pt 上回る．どこでも最も強い比較手法は 1 つも無い}
\end{frame}

\begin{frame}{基本版は 22 の組み合わせのうち \NUM{cmp-ours-vs-best-pos-n} で，最も強い比較手法を上回る}
  \rl{結果：比較の図}
  \vspace{-0.4em}
  \centering
  \begin{tikzpicture}
    \begin{groupplot}[group style={group size=1 by 2, vertical sep=0.95cm},
        width=15.2cm, height=3.3cm, ybar, /pgf/bar width=5.2pt, xmin=-0.6, xmax=11.6, ymin=-2.5, ymax=10.5,
        xtick=data, xticklabel style={font=\tiny}, ytick={-2,0,2,4,6,8,10}, ymajorgrids, grid style={mMuted!25},
        ylabel={出発点との差}, ylabel style={font=\tiny}]
      \nextgroupplot[xticklabels from table={figures/compare_k0.csv}{tick},
          title={$K=0$}, title style={font=\scriptsize, at={(0.01,0.97)}, anchor=north west}, clip=false,
          legend style={font=\tiny, draw=none, at={(1,1.12)}, anchor=south east, legend columns=3, /tikz/every even column/.append style={column sep=0.8em}}]
        \addplot[fill=cBest, draw=none, error bars/.cd, y dir=both, y explicit, error mark=none, error bar style={black, line width=0.5pt}] table[x=pos, y=b, y error plus expr=\thisrow{b_hi}-\thisrow{b}, y error minus expr=\thisrow{b}-\thisrow{b_lo}] {compare_k0.csv}; \addlegendentry{最も強い比較手法（調整版）}
        \addplot[fill=cOurs, draw=none, error bars/.cd, y dir=both, y explicit, error mark=none, error bar style={black, line width=0.5pt}] table[x=pos, y=o, y error plus expr=\thisrow{o_hi}-\thisrow{o}, y error minus expr=\thisrow{o}-\thisrow{o_lo}] {compare_k0.csv}; \addlegendentry{基本版}
        \addplot[fill=cExt, draw=none, error bars/.cd, y dir=both, y explicit, error mark=none, error bar style={black, line width=0.5pt}] table[x=pos, y=e, y error plus expr=\thisrow{e_hi}-\thisrow{e}, y error minus expr=\thisrow{e}-\thisrow{e_lo}] {compare_k0.csv}; \addlegendentry{拡張版}
        \draw[mMuted, densely dashed] (axis cs:3.75,-2.5) -- (axis cs:3.75,10.5);
        \draw[mMuted, densely dashed] (axis cs:7.25,-2.5) -- (axis cs:7.25,10.5);
        \node[font=\tiny, text=mMuted, anchor=south] at (axis cs:1.5,10.5) {1 つ目のエンコーダ，AVT の世界};
        \node[font=\tiny, text=mMuted, anchor=south] at (axis cs:5.5,10.5) {融合の生徒};
        \node[font=\tiny, text=mMuted, anchor=south] at (axis cs:9.5,10.5) {2 つ目のエンコーダ，AVT の世界};
      \nextgroupplot[xticklabels from table={figures/compare_k5.csv}{tick},
          title={$K=5$}, title style={font=\scriptsize, at={(0.01,0.97)}, anchor=north west}]
        \addplot[fill=cBest, draw=none, error bars/.cd, y dir=both, y explicit, error mark=none, error bar style={black, line width=0.5pt}] table[x=pos, y=b, y error plus expr=\thisrow{b_hi}-\thisrow{b}, y error minus expr=\thisrow{b}-\thisrow{b_lo}] {compare_k5.csv};
        \addplot[fill=cOurs, draw=none, error bars/.cd, y dir=both, y explicit, error mark=none, error bar style={black, line width=0.5pt}] table[x=pos, y=o, y error plus expr=\thisrow{o_hi}-\thisrow{o}, y error minus expr=\thisrow{o}-\thisrow{o_lo}] {compare_k5.csv};
        \addplot[fill=cExt, draw=none, error bars/.cd, y dir=both, y explicit, error mark=none, error bar style={black, line width=0.5pt}] table[x=pos, y=e, y error plus expr=\thisrow{e_hi}-\thisrow{e}, y error minus expr=\thisrow{e}-\thisrow{e_lo}] {compare_k5.csv};
        \draw[mMuted, densely dashed] (axis cs:3.75,-2.5) -- (axis cs:3.75,10.5);
        \draw[mMuted, densely dashed] (axis cs:7.25,-2.5) -- (axis cs:7.25,10.5);
    \end{groupplot}
  \end{tikzpicture}\par
  \vspace{0.1em}
  {\raggedright\tiny 目盛りの括弧はその生徒と $K$ で最も強い比較手法．棒はターゲットの平均，ひげはターゲットの間の最小から最大（95\% 区間ではない）．差は UAR のポイント\par}
  {\raggedright\scriptsize 基本版が届かないのは $K=5$ の AVT の音声と言葉（一般化蒸留に負ける）と AT の融合（共訓練と同じ）\par}
\end{frame}

\begin{frame}{最も強い比較手法との差は，同じ分割の反復をまとめても多くの生徒で有意である}
  \rl{結果：検定}
  \vspace{0.3em}
  \begin{columns}[T,onlytextwidth]
    \begin{column}{0.5\textwidth}
      \centering\footnotesize
      \renewcommand{\arraystretch}{1.3}
      \begin{tabular}{@{}lcc@{}}
        \toprule
        拡張版 − 最も強い比較手法 & $K=0$ & $K=5$ \\
        \midrule
        差が正の生徒 & \NUM{st-a-k0-pos-n} / \NUM{st-a-k0-cells-n} & \NUM{st-a-k5-pos-n} / \NUM{st-a-k5-cells-n} \\
        効果量 $d_z$ & \NUM{st-a-k0-dz-lo}〜\NUM{st-a-k0-dz-hi} & \NUM{st-a-k5-dz-lo}〜\NUM{st-a-k5-dz-hi} \\
        5\% で有意な生徒 & \NUM{st-a-k0-sig-n} & \NUM{st-a-k5-sig-n} \\
        同じ分割の反復を平均した後 & \NUM{st-a-k0-sigdedup-n} & \NUM{st-a-k5-sigdedup-n} \\
        \bottomrule
      \end{tabular}
    \end{column}
    \begin{column}{0.48\textwidth}
      \begin{itemize}\scriptsize\setlength{\itemsep}{5pt}
        \item 両側の Wilcoxon 符号順位検定を，ターゲットと反復の対で行う．Holm の補正（たくさん検定した分だけ基準を厳しくする手続き）は全部の生徒と $K$ にかける
        \item $d_z$：差の平均を差の標準偏差で割った値
        \item 平均した後に有意なのは，$K=0$ は AVT の音声の生徒と VT・AT の融合の生徒．$K=5$ は AVT の顔と言葉の生徒，AV と AT の融合の生徒を除く全部
        \item ターゲットは世界ごとに 4〜6 しかないので，この検定は私たちのコーパスの中のばらつきを表し，ほかのコーパスへの一般化は表さない
      \end{itemize}
    \end{column}
  \end{columns}
\end{frame}

\begin{frame}{$K$ を変えても，拡張版は音声の生徒ですべての $K$ で最も強い比較手法を上回る}
  \rl{結果：$K$ ごとの比較}
  \vspace{-0.3em}
  \begin{columns}[T,onlytextwidth]
    \begin{column}{0.55\textwidth}
      \centering
      \begin{tikzpicture}
        \begin{groupplot}[group style={group size=2 by 2, horizontal sep=0.35cm, vertical sep=0.55cm, xlabels at=edge bottom, ylabels at=edge left, yticklabels at=edge left, xticklabels at=edge bottom},
            width=4.7cm, height=3.4cm, xmin=-0.3, xmax=4.3, ymin=-1.2, ymax=5.5, xtick={0,...,4}, xticklabels={0,1,5,10,25},
            xlabel={$K$}, ylabel={出発点との差}, ymajorgrids, grid style={mMuted!25}, ticklabel style={font=\tiny}, label style={font=\tiny},
            title style={font=\tiny, at={(0.03,0.97)}, anchor=north west}]
          \nextgroupplot[title={音声の生徒}, legend to name=kcleg, legend style={font=\tiny, draw=none, legend columns=3, /tikz/every even column/.append style={column sep=0.8em}}]
            \draw[black, line width=0.5pt] (axis cs:-0.3,0) -- (axis cs:4.3,0);
            \addplot[cBest, mark=square*, mark size=1.4pt, line width=1pt] table[x=x, y=b_a] {kcurves.csv}; \addlegendentry{最も強い比較手法（$K$ ごとに選び直す）}
            \addplot[cOurs, mark=*, mark size=1.4pt, line width=1pt] table[x=x, y=o_a] {kcurves.csv}; \addlegendentry{基本版}
            \addplot[cExt, mark=diamond*, mark size=1.8pt, line width=1pt] table[x=x, y=x_a] {kcurves.csv}; \addlegendentry{拡張版}
          \nextgroupplot[title={顔の生徒}]
            \draw[black, line width=0.5pt] (axis cs:-0.3,0) -- (axis cs:4.3,0);
            \addplot[cBest, mark=square*, mark size=1.4pt, line width=1pt] table[x=x, y=b_v] {kcurves.csv};
            \addplot[cOurs, mark=*, mark size=1.4pt, line width=1pt] table[x=x, y=o_v] {kcurves.csv};
            \addplot[cExt, mark=diamond*, mark size=1.8pt, line width=1pt] table[x=x, y=x_v] {kcurves.csv};
          \nextgroupplot[title={言葉の生徒}]
            \draw[black, line width=0.5pt] (axis cs:-0.3,0) -- (axis cs:4.3,0);
            \addplot[cBest, mark=square*, mark size=1.4pt, line width=1pt] table[x=x, y=b_t] {kcurves.csv};
            \addplot[cOurs, mark=*, mark size=1.4pt, line width=1pt] table[x=x, y=o_t] {kcurves.csv};
            \addplot[cExt, mark=diamond*, mark size=1.8pt, line width=1pt] table[x=x, y=x_t] {kcurves.csv};
          \nextgroupplot[title={融合の生徒}]
            \draw[black, line width=0.5pt] (axis cs:-0.3,0) -- (axis cs:4.3,0);
            \addplot[cBest, mark=square*, mark size=1.4pt, line width=1pt] table[x=x, y=b_f] {kcurves.csv};
            \addplot[cOurs, mark=*, mark size=1.4pt, line width=1pt] table[x=x, y=o_f] {kcurves.csv};
            \addplot[cExt, mark=diamond*, mark size=1.8pt, line width=1pt] table[x=x, y=x_f] {kcurves.csv};
        \end{groupplot}
      \end{tikzpicture}\par
      \pgfplotslegendfromname{kcleg}\par
      {\raggedright\tiny AVT の世界，1 つ目のエンコーダ，ターゲット 4 つ，反復 10 回．差は UAR のポイント\par}
    \end{column}
    \begin{column}{0.43\textwidth}
      \begin{itemize}\scriptsize\setlength{\itemsep}{5pt}
        \item 音声の生徒では，拡張版がすべての $K$ で最も強い比較手法を \NUM{cmp-oursx-vs-best-avt-a-lo}〜\NUM{cmp-oursx-vs-best-avt-a-hi} \pt 上回る
        \item 拡張版を上回るのは，言葉の生徒の $K\ge10$ での一般化蒸留だけ（\NUM{cmp-gd-vs-oursx-avt-t-kge10-lo}〜\NUM{cmp-gd-vs-oursx-avt-t-kge10-hi} \pt）
        \item 拡張版は出発点を，音声の生徒ですべての $K$ に \NUM{cmp-oursx-vs-e39-avt-a-lo}〜\NUM{cmp-oursx-vs-e39-avt-a-hi}，言葉の生徒で $K\le5$ に \NUM{cmp-oursx-vs-e39-avt-t-kle5-lo}〜\NUM{cmp-oursx-vs-e39-avt-t-kle5-hi}，融合の生徒で $K\ge1$ に \NUM{cmp-oursx-vs-e39-avt-f-kge1-lo}〜\NUM{cmp-oursx-vs-e39-avt-f-kge1-hi} \pt 上回る
        \item 基本版は $K\ge10$ で出発点を下回ることがあるが，拡張版は下回らない
      \end{itemize}
    \end{column}
  \end{columns}
\end{frame}

\begin{frame}{固定版のエントロピー最小化は，ターゲットと世界で大きく得も損もする}
  \rl{結果：固定版}
  \vspace{-0.2em}
  \begin{columns}[T,onlytextwidth]
    \begin{column}{0.5\textwidth}
      \centering
      \begin{tikzpicture}
        \begin{axis}[width=7.2cm, height=5.0cm, xmin=-0.5, xmax=3.5, ymin=-30, ymax=8,
            xtick={0,...,3}, xticklabels={AVT,VT,AT,AV}, xlabel={融合の生徒の世界（$K=0$）},
            ylabel={出発点との差（\pt）}, ymajorgrids, grid style={mMuted!25},
            legend style={font=\tiny, draw=none, at={(0.5,1.02)}, anchor=south, legend columns=2, /tikz/every even column/.append style={column sep=0.8em}}]
          \draw[black, line width=0.5pt] (axis cs:-0.5,0) -- (axis cs:3.5,0);
          \addplot[mMuted!60, quiver={u=0.38, v=\thisrow{ox}-\thisrow{ent}}, -] table[x expr=\thisrow{x}-0.19, y=ent] {entmin.csv};
          \addplot[only marks, mark=*, mark size=1.8pt, cEnt] table[x expr=\thisrow{x}-0.19, y=ent] {entmin.csv}; \addlegendentry{エントロピー最小化（固定版）}
          \addplot[only marks, mark=diamond*, mark size=2.2pt, cExt] table[x expr=\thisrow{x}+0.19, y=ox] {entmin.csv}; \addlegendentry{拡張版（調整版）}
          \addplot[only marks, mark=-, mark size=6pt, black, line width=1.2pt] table[x expr=\thisrow{x}-0.19, y=ent] {entmin_mean.csv};
          \addplot[only marks, mark=-, mark size=6pt, black, line width=1.2pt] table[x expr=\thisrow{x}+0.19, y=ox] {entmin_mean.csv};
        \end{axis}
      \end{tikzpicture}\par
      {\raggedright\tiny 点はターゲット，黒い横棒はターゲットの平均，灰の線は同じターゲットの 2 点を結ぶ．反復 10 回\par}
    \end{column}
    \begin{column}{0.48\textwidth}
      \begin{itemize}\scriptsize\setlength{\itemsep}{5pt}
        \item 元の論文の設定では，エントロピー最小化・共正則化・事前確率の調整が，AVT の単独の生徒を $K=0$ で 1 つの感情に潰し，出発点の UAR を \NUM{cmp-fixed-collapse-avt-mod-k0-lo}〜\NUM{cmp-fixed-collapse-avt-mod-k0-hi} \pt 変える
        \item 融合の生徒のエントロピー最小化は，4 つの世界の平均で \NUM{cmp-entmin-fixed-vs-e39-f-k0-lo}〜\NUM{cmp-entmin-fixed-vs-e39-f-k0-hi} \pt
        \item IEMOCAP 即興では AVT で $+$\NUM{cmp-entmin-fixed-vs-e39-avt-f-k0-impro}，AT で $+$\NUM{cmp-entmin-fixed-vs-e39-at-f-k0-impro} だが，同じターゲットで VT \NUM{cmp-entmin-fixed-vs-e39-vt-f-k0-impro}，AV \NUM{cmp-entmin-fixed-vs-e39-av-f-k0-impro}．AVT の MELD では \NUM{cmp-entmin-fixed-vs-e39-avt-f-k0-meld}
        \item 仮のターゲットで選ぶとこの失敗を見つけ，AVT の融合の生徒ではどの $K$ でもエントロピー最小化を使わない
      \end{itemize}
    \end{column}
  \end{columns}
\end{frame}

\begin{frame}{音声のエンコーダを微調整しても，凍結した特徴での手法に届かない}
  \rl{結果：微調整との比較}
  \vspace{-0.2em}
  \begin{columns}[T,onlytextwidth]
    \begin{column}{0.53\textwidth}
      \centering
      \begin{tikzpicture}
        \begin{groupplot}[group style={group size=2 by 1, horizontal sep=0.25cm, yticklabels at=edge left},
            height=5.0cm, ybar, ymin=25, ymax=55, ymajorgrids, grid style={mMuted!25},
            xtick={0,1,2}, xticklabels={凍結,微調整,{微調整＋中心化}}, xticklabel style={font=\tiny},
            ylabel={UAR（\%）}, ticklabel style={font=\tiny}, label style={font=\tiny},
            title style={font=\scriptsize, at={(0.03,0.97)}, anchor=north west}, unbounded coords=jump]
          \nextgroupplot[width=4.0cm, /pgf/bar width=7pt, xmin=-0.5, xmax=2.5, title={$K=0$}, ymin=25, ymax=55]
            \addplot[fill=mDarkTeal!70, draw=none] table[x=x, y=start] {finetune_k0.csv};
            \addplot[fill=cOurs, draw=none] table[x=x, y=weak] {finetune_k0.csv};
          \nextgroupplot[width=5.3cm, /pgf/bar width=6pt, xmin=-0.5, xmax=2.5, title={$K=5$}, ymin=25, ymax=55,
              legend style={font=\tiny, draw=none, at={(0.5,-0.32)}, anchor=north, legend columns=2}]
            \addplot[fill=mDarkTeal!70, draw=none] table[x=x, y=konly] {finetune_k5.csv}; \addlegendentry{出発点のまま／$K$ 本だけ}
            \addplot[fill=mMuted!45, draw=none] table[x=x, y=anchor] {finetune_k5.csv}; \addlegendentry{$K$ 本＋出発点への罰則}
            \addplot[fill=cOurs, draw=none] table[x=x, y=weak] {finetune_k5.csv}; \addlegendentry{弱いラベルを足す}
        \end{groupplot}
      \end{tikzpicture}\par
      {\raggedright\tiny 音声の生徒．ターゲット 4 つ（IEMOCAP 即興・台本，MELD，CMU-MOSEI），反復 3 回．凍結の $K=5$ の「$K$ 本＋罰則」は 5 つの値を選んだ出発点．棒は偶然の水準 25\% から立てた\par}
    \end{column}
    \begin{column}{0.45\textwidth}
      \begin{itemize}\scriptsize\setlength{\itemsep}{5pt}
        \item 凍結した特徴の設定が，微調整でできることを小さく見せていないかを確かめた（WavLM を低ランクの適応器 LoRA で微調整）
        \item $K=5$ 本だけで微調整すると，凍結した 5 つの値の出発点を \NUM{cmp-ft-konly-below-frozen-e39-k5} \pt 下回る（ターゲットごとに \NUM{cmp-ft-konly-below-frozen-e39-k5-tgt-lo}〜\NUM{cmp-ft-konly-below-frozen-e39-k5-tgt-hi}）
        \item 弱いラベルは微調整したモデルも上げる：$K=0$ で \NUM{cmp-ft-weak-vs-start-k0}，$K=5$ で \NUM{cmp-ft-weak-vs-anchor-k5} \pt
        \item それでも微調整版の手法は，凍結した特徴の同じ手法より \NUM{cmp-ft-ours-below-frozen-ours-k0}（$K=0$），\NUM{cmp-ft-ours-below-frozen-ours-k5}（$K=5$）\pt 低い．中心化は差の \NUM{cmp-ftc-gap-closed-start-k0}〜\NUM{cmp-ftc-gap-closed-konly-k5}\% しか埋めない
      \end{itemize}
    \end{column}
  \end{columns}
\end{frame}

\begin{frame}{2 つ目のエンコーダでも結論は同じで，仮のターゲットで選ぶ損は手法を有利にしない}
  \rl{結果：2 つ目のエンコーダと選び方}
  \vspace{0.3em}
  \begin{columns}[T,onlytextwidth]
    \begin{column}{0.49\textwidth}
      \textbf{\footnotesize 2 つ目のエンコーダ（HuBERT，CLIP，DeBERTa）}\par\vspace{0.2em}
      \begin{itemize}\scriptsize\setlength{\itemsep}{5pt}
        \item 基本版は $K=0$ と $K=5$ で，4 つの生徒すべてを出発点より上げる
        \item 拡張版は最も強い比較手法を，$K=0$・5 で \NUM{cmp-oursx-vs-best-w2-k0k5-lo}〜\NUM{cmp-oursx-vs-best-w2-k0k5-hi} \pt，$K=25$ で \NUM{cmp-oursx-vs-best-w2-k25-lo}〜\NUM{cmp-oursx-vs-best-w2-k25-hi} \pt 上回る
      \end{itemize}
    \end{column}
    \begin{column}{0.49\textwidth}
      \textbf{\footnotesize 選び方で失う分（$K=0$）}\par\vspace{0.2em}
      \begin{itemize}\scriptsize\setlength{\itemsep}{5pt}
        \item テスト話者で最良の候補（オラクル）は，選んだ候補を AVT の言葉の生徒で \NUM{cmp-oursx-oracle-vs-tuned-avt-t-k0}，2 つ目のエンコーダの言葉の生徒で \NUM{cmp-oursx-oracle-vs-tuned-w2-t-k0}，ほかの生徒で \NUM{cmp-oursx-oracle-vs-tuned-at-f-k0}〜\NUM{cmp-oursx-oracle-vs-tuned-w2-v-k0} \pt 上回る
        \item 差が 1 \pt に届くのは \NUM{cmp-oursx-oracle-gap-students-n} の生徒のうち \NUM{cmp-oursx-oracle-gap-ge1-n}．よりよい選び方で足せる分の上限になる
        \item 比較手法も同じ手順で選ぶので，この損は私たちの手法を有利にしない
      \end{itemize}
    \end{column}
  \end{columns}
  \vspace{0.6em}
  \takeaway{オラクルの値は選び方の損を測るためだけに使い，手法の性能としては出さない}
\end{frame}
